\chapter{Sprint 5 — Search Orchestrator, Voice Transcription and Portals}

\section*{Introduction}
\addcontentsline{toc}{section}{Introduction}

Sprint 5 builds the interactive surface of the BVIRAL Platform: an \textbf{Agentic Helper}
service for chatbot and video discovery, a deterministic \textbf{v2 search orchestrator} over
the hybrid search engine, a structured \textbf{filter extractor} that turns free text into typed
parameters, an \textbf{OpenAI Whisper} voice transcription pipeline, and the two \textbf{Next.js
portals} that wire the end-to-end experience together through BFF routes. Normal video search now
uses a faster typed plan plus one advanced-search call.

%% ─── Sprint Backlog ─────────────────────────────────────────────────────────
\section{Sprint Backlog}

\begin{table}[H]
\centering
\caption{Sprint 5 Backlog}
\label{tab:sprint5_backlog}
\begin{tabularx}{\textwidth}{clXcc}
\toprule
\textbf{ID} & \textbf{Story} & \textbf{Task} & \textbf{Effort (d)} & \textbf{Prio.} \\
\midrule
T5.1 & US-11 & \texttt{VideoSearchFilterPlan} structured LLM output              & 1 & M \\
T5.2 & US-11 & Deterministic \texttt{VideoSearchService} over \texttt{/advanced-search} & 1 & M \\
T5.3 & US-11 & ReAct agent retained as fallback for complex tool workflows & 1 & M \\
T5.4 & US-11 & Search state, pagination follow-ups, and response metadata             & 0.5 & S \\
T5.5 & US-13 & \texttt{/api/v1/transcribe} endpoint forwarding to Whisper        & 1 & M \\
T5.6 & US-13 & Next.js \texttt{/api/transcribe} BFF route (proxy)                 & 0.5 & M \\
T5.7 & US-13 & Mic button with \texttt{MediaRecorder} in both portals             & 1 & M \\
T5.8 & —     & Video Search Portal (Next.js 15, Radix UI, Tailwind)              & 2 & M \\
T5.9 & —     & Authentic Rights Portal (Next.js 16, Material UI, Stripe.js)       & 2 & M \\
T5.10 & —     & ChatbotWidget component shared by both portals                    & 1 & M \\
T5.11 & US-11 & Recommendation events API (\texttt{/recommendations/events/*})      & 1 & M \\
T5.12 & US-11 & User interest profile model + blend ranking service               & 1 & M \\
T5.13 & US-11 & Search portal instrumentation (search/click events)               & 1 & M \\
T5.14 & US-11 & Anti-repetition ranking (seen suppression, freshness, exploration) & 1 & M \\
T5.15 & US-11 & Query rewriting for retrieval-optimised RAG queries               & 0.5 & S \\
T5.16 & US-11 & Relevance grading with corrective re-retrieval (CRAG pattern)     & 1   & S \\
T5.17 & US-11 & Semantic query enrichment for video search (vector optimisation)  & 0.5 & S \\
\bottomrule
\end{tabularx}
\end{table}

%% ─── Design ─────────────────────────────────────────────────────────────────
\section{Design}

\subsection{Use Case Diagram for Sprint 5}

Sprint 5 surfaces the \emph{Invoke video-search agent} and \emph{Transcribe audio query} use cases
presented in the global use-case diagram (Figure~\ref{fig:uc_global}) and in
Figure~\ref{fig:uc_chat}.

\subsection{Class Diagram for Sprint 5}

The first implementation used an in-memory conversation window and a ReAct loop. The
v2 architecture keeps the same public API but introduces a typed orchestration layer: a unified
router classifies each request, Qdrant powers RAG answers, Redis can persist short-term memory
with TTL, and \texttt{VideoSearchService} converts free-text search requests into a validated
\texttt{VideoSearchPlan}. The key Sprint 5 model remains the structured search plan, now aligned
with the advanced hybrid search endpoint.

\begin{lstlisting}[language=python, caption={VideoSearchFilterPlan Pydantic model}]
class VideoSearchFilterPlan(BaseModel):
    method: Literal["auto", "semantic", "filters"] = "auto"
    query: str = ""
    k: int = Field(default=10, ge=1, le=50)
    categories: list[str] = Field(default_factory=list)
    tags: list[str] = Field(default_factory=list)
    duration_min: float | None = None
    duration_max: float | None = None
    locations: list[str] = Field(default_factory=list)
    resolutions: list[str] = Field(default_factory=list)
    orientation: list[str] = Field(default_factory=list)
    video_id: str | None = None
    offset: int = Field(default=0, ge=0)
    limit: int = Field(default=10, ge=1, le=50)
\end{lstlisting}

\subsection{Sequence Diagrams}

\subsubsection{v2 Auto Routing and Advanced Search}

Figure~\ref{fig:seq_agent} shows the current v2 flow. A request first goes through the
structured router. If the route is \texttt{VIDEO\_SEARCH}, the service loads the previous search
state, extracts a typed plan, normalizes it deterministically, and calls the advanced hybrid
search endpoint once. ReAct is not on the normal path; it remains a fallback for future multi-tool
exploration.

\begin{figure}[H]
\centering
\begin{tikzpicture}[x=1cm, y=1cm, font=\fontsize{7}{8}\selectfont]
\node[sdFE]   (pUser)   at ( 0.0, 0) {User};
\node[sdFE]   (pBFF)    at ( 2.5, 0) {Portal BFF};
\node[sdBE]   (pChat)   at ( 5.0, 0) {Chat API};
\node[sdExt]  (pRouter) at ( 7.5, 0) {Router LLM};
\node[sdData] (pMem)    at (10.0, 0) {Redis\\Memory};
\node[sdBE]   (pSearch) at (12.5, 0) {VideoSearch\\Service};
\node[sdData] (pVideo)  at (15.0, 0) {Videos API\\Advanced};

\foreach \x in {0.0, 2.5, 5.0, 7.5, 10.0, 12.5, 15.0}
  \draw[dashed, gray!50] (\x,-0.3) -- (\x,-12.2);

\draw[sdmsg] (0.0,-0.85) -- node[above, pos=0.5]{search request} (2.5,-0.85);
\draw[sdmsg] (2.5,-1.55) -- node[above, pos=0.5]{POST /api/v1/chat} (5.0,-1.55);
\draw[sdmsg] (5.0,-2.25) -- node[above, pos=0.5]{classify route} (7.5,-2.25);
\draw[sdret] (7.5,-2.95) -- node[above, pos=0.5]{VIDEO\_SEARCH + confidence} (5.0,-2.95);
\draw[sdmsg] (5.0,-3.65) -- node[above, pos=0.5]{load history + last plan} (10.0,-3.65);
\draw[sdret] (10.0,-4.35) -- node[above, pos=0.5]{window + search state} (5.0,-4.35);
\draw[sdmsg] (5.0,-5.05) -- node[above, pos=0.5]{question, history, previous plan} (12.5,-5.05);
\draw[sdmsg] (12.5,-5.75) to[out=-30,in=30] node[right, pos=0.5]{extract typed plan} (12.5,-6.45);
\draw[sdmsg] (12.5,-7.15) to[out=-30,in=30] node[right, pos=0.5]{normalize filters + pagination} (12.5,-7.85);
\draw[sdmsg] (12.5,-8.55) -- node[above, pos=0.5]{POST /advanced-search} (15.0,-8.55);
\draw[sdret] (15.0,-9.25) -- node[above, pos=0.5]{items, total, execution} (12.5,-9.25);
\draw[sdmsg] (12.5,-9.95) -- node[above, pos=0.5]{save last plan} (10.0,-9.95);
\draw[sdret] (12.5,-10.65) -- node[above, pos=0.5]{videos + search action} (5.0,-10.65);
\draw[sdret] (5.0,-11.35) -- node[above, pos=0.5]{ChatResponse} (0.0,-11.35);

\node[sdnote] at (14.8,-6.35) {Fallback path\\only if advanced\\search is disabled};
\draw[->, dashed, gray!55] (12.5,-6.45) -- (14.1,-6.35);

\end{tikzpicture}
\caption{Agentic Helper v2 video search sequence — router, typed plan and advanced search}
\label{fig:seq_agent}
\end{figure}

\subsubsection{Audio Transcription}

Figure~\ref{fig:seq_transcribe} captures the voice-to-chat pipeline, from the browser mic button
to the Agentic Helper proxy to OpenAI Whisper.

\begin{figure}[H]
\centering
\begin{tikzpicture}[x=1cm, y=1cm, font=\fontsize{7}{8}\selectfont]
\node[sdFE]  (pUser)    at (0,  0) {User\\(browser)};
\node[sdFE]  (pPortal)  at (3,  0) {Portal BFF\\(Next.js)};
\node[sdBE]  (pAgent)   at (6,  0) {Agentic\\Helper};
\node[sdExt] (pWhisper) at (9,  0) {OpenAI\\Whisper};

\foreach \x in {0, 3, 6, 9}
  \draw[dashed, gray!50] (\x,-0.3) -- (\x,-9.5);

\draw[sdmsg] (0,-0.85) -- node[above, pos=0.5]{click mic button} (3,-0.85);
\draw[sdmsg] (3,-1.57) -- node[above, pos=0.5]{MediaRecorder.start()} (0,-1.57);
\draw[sdmsg] (0,-2.29) to[out=-30,in=30] node[right, pos=0.5]{stream audio chunks} (0,-3.01);
\draw[sdmsg] (0,-3.73) -- node[above, pos=0.5]{stop recording (audio Blob)} (3,-3.73);
\draw[sdmsg] (3,-4.45) -- node[above, pos=0.5]{POST /api/transcribe (multipart)} (6,-4.45);
\draw[sdmsg] (6,-5.17) -- node[above, pos=0.5]{POST /v1/audio/transcriptions (Whisper)} (9,-5.17);
\draw[sdret] (9,-5.89) -- node[above, pos=0.5]{\{text: "transcribed content"\}} (6,-5.89);
\draw[sdret] (6,-6.61) -- node[above, pos=0.5]{\{transcription\}} (3,-6.61);
\draw[sdret] (3,-7.33) -- node[above, pos=0.5]{pre-fill chat input} (0,-7.33);
\draw[sdmsg] (0,-8.05) -- node[above, pos=0.5]{(optional) edit + submit as message} (3,-8.05);

\node[sdnote] at (11.0,-5.53) {Supports mp3,\\mp4, wav, webm\\up to 25MB};
\draw[->, dashed, gray!50] (9,-5.53) -- (10.3,-5.53);

\end{tikzpicture}
\caption{Audio transcription sequence diagram — Whisper pipeline}
\label{fig:seq_transcribe}
\end{figure}

%% ─── Implementation ────────────────────────────────────────────────────────
\section{Implementation}

\subsection{v2 Agentic Helper Architecture}

The optimized architecture treats the chatbot as an orchestrated backend system instead of a
single prompt. The \texttt{/api/v1/chat} endpoint accepts three modes:
\texttt{default}, \texttt{video\_search}, and \texttt{auto}. In \texttt{auto} mode, a structured
router returns one of five routes: \texttt{CHAT\_RAG}, \texttt{VIDEO\_SEARCH},
\texttt{DIRECT}, \texttt{OUT\_OF\_SCOPE}, or \texttt{SUPPORT\_HANDOFF}. This makes routing
observable and testable before any expensive retrieval or tool execution begins.

\begin{figure}[H]
\centering
\begin{tikzpicture}[
  x=1cm, y=1cm,
  block/.style={draw=teal!70!black, fill=teal!8, rounded corners=4pt,
    minimum width=3.2cm, minimum height=0.78cm, align=center,
    font=\fontsize{8}{9}\selectfont},
  data/.style={draw=violet!70!black, fill=violet!8, rounded corners=4pt,
    minimum width=3.0cm, minimum height=0.78cm, align=center,
    font=\fontsize{8}{9}\selectfont},
  ext/.style={draw=gray!65, fill=gray!8, rounded corners=4pt,
    minimum width=3.0cm, minimum height=0.78cm, align=center,
    font=\fontsize{8}{9}\selectfont},
  arr/.style={->, >=Stealth, thick}
]
\node[block] (api) at (0,0) {\texttt{/api/v1/chat}\\FastAPI endpoint};
\node[block] (router) at (0,-1.35) {Structured router\\route + confidence};
\node[block] (rag) at (-4.0,-3.0) {RAG answer chain\\policy + citations};
\node[block] (search) at (0,-3.0) {VideoSearchService\\typed plan executor};
\node[block] (direct) at (4.3,-3.0) {Direct / handoff\\guarded response};
\node[data] (qdrant) at (-4.0,-4.55) {Qdrant\\knowledge chunks};
\node[ext] (hybrid) at (0,-4.55) {Videos Search API\\\texttt{/advanced-search}};
\node[data] (redis) at (4.3,-4.55) {Redis memory\\window + state TTL};
\node[ext] (fallback) at (0,-7.25) {Fallback path\\optional multi-tool path};
\node[block] (resp) at (0,-6.0) {ChatResponse\\answer, videos, action, metadata};

\draw[arr] (api) -- (router);
\draw[arr] (router) -- (rag);
\draw[arr] (router) -- (search);
\draw[arr] (router) -- (direct);
\draw[arr] (rag) -- (qdrant);
\draw[arr] (search) -- (hybrid);
\draw[arr] (router) -- (redis);
\draw[arr, dashed, gray!70] (search) -- (fallback);
\draw[arr] (rag) |- (resp);
\draw[arr] (search) -- (resp);
\draw[arr] (direct) |- (resp);
\end{tikzpicture}
\caption{Agentic Helper v2 — explicit routing, deterministic search, RAG and memory}
\label{fig:agentic_helper_v2}
\end{figure}

\subsection{Deterministic Advanced Video Search}

The original ReAct loop was useful for prototyping, but it was not the best default for a
high-volume search path. The optimized path is deterministic:

\begin{enumerate}
  \item \texttt{VideoSearchPlan} is extracted with structured output.
  \item Server-side normalization applies hard rules for top/trending, newest/oldest,
    shortest/longest, exact video IDs, categories, orientation, duration, resolution, and
    pagination follow-ups.
  \item The plan is converted to the \texttt{/api/videos/advanced-search} payload:
    \texttt{query}, \texttt{filters}, \texttt{sort}, \texttt{pagination}, and optional
    \texttt{debug}.
  \item The client uses strict timeouts, retries, API-key headers, and base-URL fallback.
  \item The response is shaped into video cards, \texttt{search\_filters},
    \texttt{search\_action}, \texttt{search\_url}, \texttt{total}, \texttt{next\_offset},
    \texttt{execution}, and \texttt{fallbacks\_used}.
\end{enumerate}

This design uses the advanced search engine from Sprint 3 as the single ranking source of truth.
It avoids unnecessary agent loops for normal searches, reduces token cost, and makes search
behavior easier to test with deterministic unit tests. The previous ReAct agent remains behind the
service boundary as a fallback for future exploratory workflows that genuinely need multiple tool
calls.

\subsection{Semantic Query Enrichment for Video Search}

The deterministic plan extractor produces a \texttt{plan.query} that is the semantic portion of
the user request after filters have been stripped. While this is already an improvement over
passing the raw message to OpenSearch, it still has a vocabulary mismatch issue: a short,
conversational phrase such as \emph{``funny dog videos''} may not align closely with the dense
embedding space learned from video titles, descriptions, and tags.

To close this gap, an optional \textbf{semantic query enrichment} step runs between
\texttt{extract\_plan} and \texttt{normalize\_plan} inside \texttt{VideoSearchService.search()}:

\begin{enumerate}
  \item The filter extractor produces the raw \texttt{VideoSearchPlan}, including a terse
    \texttt{plan.query}.
  \item If \texttt{plan.query} is non-empty (i.e., the request is not a ranking-only query),
    \texttt{LLMService.rewrite\_video\_search\_query()} expands it into a richer, more
    descriptive form. The classifier-grade LLM is instructed to add synonyms, visual
    descriptors, and contextual terms likely to appear in matching video metadata, while keeping
    the result under 20 words.
  \item The enriched query replaces \texttt{plan.query} before normalization and before the
    \texttt{/api/videos/advanced-search} call.
\end{enumerate}

Ranking-only requests (where the filter extractor correctly leaves \texttt{query=""} because
the user asked for ``top trending'' or ``most viewed'') are automatically skipped — no
unnecessary LLM call is made. The feature is gated by
\texttt{video\_search\_query\_rewriting\_enabled} so it can be toggled in \texttt{.env}
independently of the RAG rewriting flag.

Figure~\ref{fig:video_query_rewrite} shows the updated \texttt{VideoSearchService} pipeline
with the enrichment step inserted.

\begin{figure}[H]
\centering
\begin{tikzpicture}[
  x=1cm, y=1cm,
  block/.style={draw=teal!70!black, fill=teal!8, rounded corners=4pt,
    minimum width=4.6cm, minimum height=0.82cm, align=center,
    font=\fontsize{8}{9}\selectfont},
  data/.style={draw=violet!70!black, fill=violet!8, rounded corners=4pt,
    minimum width=4.6cm, minimum height=0.82cm, align=center,
    font=\fontsize{8}{9}\selectfont},
  ext/.style={draw=gray!65, fill=gray!8, rounded corners=4pt,
    minimum width=3.2cm, minimum height=0.82cm, align=center,
    font=\fontsize{8}{9}\selectfont},
  new/.style={draw=orange!70!black, fill=orange!8, rounded corners=4pt,
    minimum width=4.6cm, minimum height=0.82cm, align=center,
    font=\fontsize{8}{9}\selectfont},
  dec/.style={draw=orange!70!black, fill=orange!10, diamond,
    aspect=2.6, align=center, inner sep=2pt,
    font=\fontsize{7.5}{8.5}\selectfont},
  arr/.style={->, >=Stealth, thick}
]
\node[ext]   (q)      at (0,    0)   {User question};
\node[block] (ext)    at (0,  -1.6)  {Plan Extraction\\[-1pt]\scriptsize\texttt{extract\_plan} → \texttt{VideoSearchPlan}};
\node[dec]   (chk)    at (0,  -3.3)  {query non-\\empty?};
\node[new]   (rw)     at (-3.5,-4.8) {Query Enrichment\\[-1pt]\scriptsize\texttt{rewrite\_video\_search\_query}};
\node[block] (norm)   at (0,  -6.3)  {Plan Normalisation\\[-1pt]\scriptsize\texttt{normalize\_plan}};
\node[data]  (api)    at (0,  -7.9)  {Advanced Search API\\[-1pt]\scriptsize\texttt{POST /api/videos/advanced-search}};
\node[block] (resp)   at (0,  -9.5)  {VideoSearchResult\\[-1pt]\scriptsize videos, filters, answer};

\draw[arr] (q)    -- (ext);
\draw[arr] (ext)  -- (chk);
\draw[arr] (chk.west) -| node[above, font=\scriptsize, pos=0.2]{yes} (rw.north);
\draw[arr] (chk.south) -- node[right, font=\scriptsize]{no (skip)} (norm.north);
\draw[arr] (rw.south) |- (norm.west);
\draw[arr] (norm) -- (api);
\draw[arr] (api)  -- (resp);

\node[font=\fontsize{6.5}{8}\selectfont, text=orange!60!black, align=center]
  at (3.8,-4.8) {New step\\(T5.17)};
\end{tikzpicture}
\caption{VideoSearchService pipeline with semantic query enrichment (T5.17)}
\label{fig:video_query_rewrite}
\end{figure}

\paragraph{Example.}
For the query \emph{``funny dog videos''}, the enrichment step produces a richer form such as
\emph{``hilarious dogs playing tricks humor comedy canine pets funny clips''}, which achieves
higher cosine similarity against video embeddings derived from matching titles and descriptions.
The semantic query passes through \texttt{normalize\_plan} unchanged (normalization only touches
filters, sort, and pagination) before being sent to OpenSearch.

\subsection{Memory, Citations and Reliability}

Conversation memory now has two implementations. Local development can use
\texttt{InMemoryChatHistoryStore}, while production can configure \texttt{CHAT\_HISTORY\_REDIS\_URL}
to store the latest conversation window and search state with a TTL. Search follow-ups such as
``show more'' or ``same but vertical'' reuse the previous \texttt{VideoSearchPlan} from this
state store.

RAG responses include source metadata in the API response through \texttt{sources[]} and
\texttt{citations[]}. Search responses include execution metadata and fallback information from
the advanced search service. These fields make the assistant easier to debug without exposing
internal prompts to the frontend.

\subsection{Agentic RAG Pipeline}

The original RAG path retrieved chunks directly from the raw conversational question and passed
all returned chunks to the LLM without any relevance check. Two weaknesses motivated an upgrade:

\begin{itemize}
  \item \textbf{Vocabulary mismatch}: a conversational question such as \emph{``what about my
    cut?''} contains pronouns and implicit references that map poorly to the keyword-dense
    representations stored in Qdrant.
  \item \textbf{Context pollution}: low-scoring, off-topic chunks that pass the top-$k$
    threshold are sent to the answer model unchanged, diluting or contradicting the grounded
    evidence.
\end{itemize}

The upgraded pipeline adds two corrective steps that transform the single-hop RAG into a
multi-step \textbf{Agentic RAG} loop, drawing on the Corrective RAG (CRAG)
pattern~\cite{crag_2024} and query-rewriting techniques for
retrieval~\cite{query_rewriting_rag}.

\subsubsection{Step 1 — Query Rewriting}

Before calling the Qdrant retriever, a fast \textbf{classifier-grade LLM call}
(\texttt{LLMService.rewrite\_rag\_query}) rewrites the conversational question into a
concise, coreference-resolved, keyword-dense retrieval query. The conversation history window
is included so the model can resolve pronouns and context fragments:

\begin{center}
\textit{``what about my cut?''} \quad $\longrightarrow$ \quad
\textit{``creator revenue share percentage BVIRAL licensing agreement''}
\end{center}

The original question is preserved for the final answer chain; only the retrieval step uses the
rewritten query. The cheaper classifier model handles this step, keeping added latency minimal.

\subsubsection{Step 2 — Relevance Grading and Corrective Re-retrieval}

After retrieval, each chunk is graded independently and \textbf{in parallel}
(\texttt{LLMService.grade\_chunks}) using a structured output call
(\texttt{with\_structured\_output(ChunkGrade)}). The model assigns a binary
\texttt{relevant: true/false} per chunk. If the number of relevant chunks falls below a
configurable threshold (\texttt{rag\_min\_relevant\_chunks}, default~2), the service performs a
\textbf{corrective re-retrieval} on the original unmodified question, merges the two result
sets, deduplicates by chunk~id, and caps the combined list at \texttt{rag\_top\_k}.

The fallback re-retrieval activates only when the rewritten query under-performs — adding at
most one extra retrieval round-trip precisely on the cases that matter most. Normal requests pay
only the cost of parallel grading calls, which complete concurrently with low overhead thanks to
\texttt{asyncio.gather}.

Both steps are individually gated by feature flags
(\texttt{rag\_query\_rewriting\_enabled}, \texttt{rag\_relevance\_grading\_enabled}) so either
can be disabled in \texttt{.env} without code changes.

Figure~\ref{fig:agentic_rag} illustrates the complete pipeline.

\begin{figure}[H]
\centering
\begin{tikzpicture}[
  x=1cm, y=1cm,
  block/.style={draw=teal!70!black, fill=teal!8, rounded corners=4pt,
    minimum width=4.4cm, minimum height=0.84cm, align=center,
    font=\fontsize{8}{9}\selectfont},
  data/.style={draw=violet!70!black, fill=violet!8, rounded corners=4pt,
    minimum width=4.4cm, minimum height=0.84cm, align=center,
    font=\fontsize{8}{9}\selectfont},
  ext/.style={draw=gray!65, fill=gray!8, rounded corners=4pt,
    minimum width=3.2cm, minimum height=0.84cm, align=center,
    font=\fontsize{8}{9}\selectfont},
  dec/.style={draw=orange!70!black, fill=orange!10, diamond,
    aspect=2.8, align=center, inner sep=3pt,
    font=\fontsize{7.5}{8.5}\selectfont},
  arr/.style={->, >=Stealth, thick}
]
%% ─── Main vertical flow ────────────────────────────────────────────────────
\node[ext]   (q)     at (0,    0)   {User question};
\node[block] (rw)    at (0,  -1.7)  {Query Rewriting\\[-1pt]\scriptsize\texttt{rewrite\_rag\_query}};
\node[data]  (ret)   at (0,  -3.4)  {Qdrant Retrieval\\[-1pt]\scriptsize top-$k$ chunks};
\node[block] (grade) at (0,  -5.1)  {Relevance Grading\\[-1pt]\scriptsize\texttt{grade\_chunks} (parallel)};
\node[dec]   (chk)   at (0,  -6.9)  {relevant chunks\\$\geq$ threshold?};
\node[block] (ans)   at (0,  -9.6)  {Answer with graded context\\[-1pt]\scriptsize\texttt{ask\_with\_context}};
%% ─── Corrective fallback branch ────────────────────────────────────────────
\node[data]  (retry) at (5.8, -8.0) {Re-retrieve\\[-1pt]\scriptsize on original question};
\node[block] (merge) at (5.8, -9.6) {Merge + deduplicate\\[-1pt]\scriptsize top-$k$};
%% ─── Arrows ─────────────────────────────────────────────────────────────────
\draw[arr] (q)     -- (rw);
\draw[arr] (rw)    -- node[right, font=\scriptsize]{\textit{rewritten query}} (ret);
\draw[arr] (ret)   -- node[right, font=\scriptsize]{\textit{top-$k$ chunks}} (grade);
\draw[arr] (grade) -- (chk);
\draw[arr] (chk.south)  -- node[left, font=\scriptsize]{yes} (ans.north);
\draw[arr] (chk.east)   -| node[above, font=\scriptsize, pos=0.18]{no} (retry);
\draw[arr] (retry)  -- (merge);
\draw[arr] (merge.west) -- (ans.east);
\end{tikzpicture}
\caption{Agentic RAG pipeline — query rewriting, parallel relevance grading, and corrective re-retrieval}
\label{fig:agentic_rag}
\end{figure}

\subsubsection{Configuration}

\begin{table}[H]
\centering
\caption{Agentic RAG configuration knobs}
\label{tab:agentic_rag_config}
\begin{tabularx}{\textwidth}{llX}
\toprule
\textbf{Setting} & \textbf{Default} & \textbf{Effect} \\
\midrule
\texttt{rag\_query\_rewriting\_enabled}  & \texttt{true}  & Enable/disable Step~1 query rewriting. \\
\texttt{rag\_relevance\_grading\_enabled}& \texttt{true}  & Enable/disable Step~2 chunk grading. \\
\texttt{rag\_min\_relevant\_chunks}      & \texttt{2}     & Threshold below which corrective re-retrieval is triggered. \\
\texttt{rag\_top\_k}                     & \texttt{4}     & Maximum chunks kept after merging. \\
\bottomrule
\end{tabularx}
\end{table}

\subsection{Voice Transcription Pipeline}

The \texttt{POST /api/v1/transcribe} endpoint accepts a multipart audio file and forwards it to
\textbf{OpenAI Whisper} (\texttt{/v1/audio/transcriptions}). Both portals expose a microphone
button that:

\begin{enumerate}
  \item Starts the browser's \texttt{MediaRecorder} API and streams audio chunks.
  \item Assembles the recording into a \texttt{Blob} when the user stops speaking.
  \item Sends the blob to the Next.js BFF route (\texttt{/api/transcribe}), which proxies it
    to the Agentic Helper.
  \item Receives the transcribed text and pre-fills the chat input field.
\end{enumerate}

\subsection{Next.js Portals}

\subsubsection{Video Search Portal (\texttt{videosearchv0/})}

\begin{itemize}
  \item \textbf{Next.js 15} with the App Router and React Server Components.
  \item \textbf{Radix UI} primitives + \textbf{Tailwind CSS} for a dark, content-first UI.
  \item \textbf{Clerk} for authentication and session management.
  \item BFF routes (\texttt{app/api/chat/route.ts}, \texttt{app/api/transcribe/route.ts}) that
    forward requests to the Agentic Helper and attach the required API key/user headers.
\end{itemize}

\subsubsection{Authentic Rights Portal (\texttt{authentic\_rights\_portal/})}

\begin{itemize}
  \item \textbf{Next.js 16} with the App Router.
  \item \textbf{Material UI (MUI)} for dashboard-style components and tables.
  \item \textbf{Clerk} authentication, \textbf{Stripe.js} embedded checkout,
    \textbf{Stripe Elements} for the digital agreement signature.
  \item BFF routes for chat, transcription, and protected AR-API calls.
\end{itemize}

\subsubsection{Admin Extensions in the Authentic Rights Portal}

The rights portal was extended with an \textbf{admin workspace} dedicated to moderation,
support supervision, and knowledge-base maintenance. Four additions were particularly important:

\begin{itemize}
  \item \textbf{Submitted-video moderation}: administrators can review user-submitted clips,
    add moderation notes, and trigger a \emph{fake accept} or reject status used for internal
    validation before any real library ingestion workflow is automated.
  \item \textbf{Global conversation supervision}: an admin-only AR-API endpoint
    (\texttt{GET /conversations/admin/all}) aggregates conversation histories across users so the
    support team can inspect titles, participants, message counts, and the latest activity.
  \item \textbf{Admin chatbot tester}: a dedicated console lets administrators switch between
    \texttt{chat} mode and \texttt{video\_search} mode, test the assistant against both
    behaviors, and inspect returned text or video cards in one interface.
  \item \textbf{FAQ management on vector storage}: the FAQ manager was reconnected to the RAG
    vector store rather than a static list. Through BFF routes, the portal now consumes
    vector-management endpoints for listing, semantic search, creation, update, and deletion.
\end{itemize}

This design keeps the browser isolated from the Agentic Helper API key while still exposing a
usable knowledge-base administration surface. The FAQ management flow now maps business FAQ
records to vector chunks, where the question, answer, visibility flag, and display order are
stored in \texttt{metadata\_json}, while semantic search is delegated to
\texttt{GET /vectors/search}.

\subsubsection{Shared ChatbotWidget}

Both portals embed a shared \texttt{ChatbotWidget} component that:
\begin{itemize}
  \item Maintains the chat message list and typing state in React state.
  \item Renders video result cards directly in the conversation flow when the Agentic Helper
    returns a \texttt{videos[]} array.
  \item Includes the mic button that triggers the Whisper transcription flow.
  \item Persists conversation continuity through the \texttt{X-User-Id} header.
\end{itemize}

\subsection{FAQ Vector Administration}

To support non-developer maintenance of the support knowledge base, the Agentic Helper now exposes
a dedicated vector-management API under \texttt{/api/v1/rag/vectors}. The endpoints allow
operators to:
\begin{itemize}
  \item list stored vectors with pagination and optional source filtering;
  \item run semantic search over the FAQ collection;
  \item create a new vector from manual content;
  \item update an existing vector, automatically re-embedding edited content;
  \item delete obsolete vectors by point id.
\end{itemize}

This administrative surface is used by the portal BFF routes to power FAQ maintenance screens
without direct access to Qdrant credentials from the frontend.

\subsection{Delivered Interfaces}

Sprint 5 groups the user-facing portal surfaces delivered across the platform: the public Video
Search Portal landing pages, the authenticated dashboard views, the embedded chatbot experience,
and the rights-portal administration screens.

\begin{figure}[H]
\centering
\includegraphics[width=\textwidth]{\detokenize{figures/videosearch/videosearch portal header.png}}
\caption{Video Search Portal — hero section with semantic search bar}
\label{fig:vsp_hero}
\end{figure}

\begin{figure}[H]
\centering
\includegraphics[width=\textwidth]{\detokenize{figures/videosearch/videosearch portal latest adition.png}}
\caption{Video Search Portal — latest additions section on the authenticated dashboard}
\label{fig:vsp_latest}
\end{figure}

\begin{figure}[H]
\centering
\includegraphics[width=\textwidth]{\detokenize{figures/videosearch/videosearch portal categorie page.png}}
\caption{Video Search Portal — browse by category page}
\label{fig:vsp_categories}
\end{figure}

\begin{figure}[H]
\centering
\begin{subfigure}[t]{\textwidth}
  \centering
  \includegraphics[width=\textwidth]{\detokenize{figures/videosearch/videosearch home page part 4.png}}
  \caption{Pricing plan and platform overview}
\end{subfigure}
\vspace{0.5em}
\begin{subfigure}[t]{\textwidth}
  \centering
  \includegraphics[width=\textwidth]{\detokenize{figures/videosearch/videosearch home page part 5.png}}
  \caption{Value proposition sections}
\end{subfigure}
\caption{Video Search Portal — public landing page (pricing and value propositions)}
\label{fig:vsp_landing_1}
\end{figure}

\begin{figure}[H]
\centering
\begin{subfigure}[t]{\textwidth}
  \centering
  \includegraphics[width=\textwidth]{\detokenize{figures/videosearch/videosearch home page part 6.png}}
  \caption{Creator testimonials and blog articles}
\end{subfigure}
\vspace{0.5em}
\begin{subfigure}[t]{\textwidth}
  \centering
  \includegraphics[width=\textwidth]{\detokenize{figures/videosearch/videosearch home page part 7.png}}
  \caption{Service offerings and getting started}
\end{subfigure}
\caption{Video Search Portal — testimonials and service overview}
\label{fig:vsp_landing_2}
\end{figure}

\begin{figure}[H]
\centering
\includegraphics[width=\textwidth]{\detokenize{figures/videosearch/videosearch home page part 8.png}}
\caption{Video Search Portal — FAQ section and footer}
\label{fig:vsp_faq_footer}
\end{figure}

\begin{figure}[H]
\centering
\includegraphics[width=0.97\textwidth]{\detokenize{figures/ar/ser dashboard.png}}
\caption{Authentic Rights Portal — dashboard overview}
\label{fig:screen_ar_dashboard}
\end{figure}

\begin{figure}[H]
\centering
\includegraphics[width=0.97\textwidth]{\detokenize{figures/ar/ar manage subscription page.png}}
\caption{Authentic Rights Portal — manage subscription}
\label{fig:screen_ar_manage_subscription}
\end{figure}

\begin{figure}[H]
\centering
\includegraphics[width=0.97\textwidth]{\detokenize{figures/ar/ar support and help page.png}}
\caption{Authentic Rights Portal — support and help}
\label{fig:screen_ar_support_help}
\end{figure}

\begin{figure}[H]
\centering
\includegraphics[width=0.97\textwidth]{\detokenize{figures/ar/ar submit video page.png}}
\caption{Authentic Rights Portal — submit video form}
\label{fig:screen_ar_submit_video}
\end{figure}

\begin{figure}[H]
\centering
\includegraphics[width=0.97\textwidth]{\detokenize{figures/ar/ar playlist videos page.png}}
\caption{Authentic Rights Portal — playlist videos}
\label{fig:screen_ar_playlists}
\end{figure}

\begin{figure}[H]
\centering
\includegraphics[width=0.97\textwidth]{\detokenize{figures/ar/playlist video grid.png}}
\caption{Authentic Rights Portal — playlist video grid}
\label{fig:screen_ar_playlist_video_grid}
\end{figure}

\begin{figure}[H]
\centering
\includegraphics[width=0.55\textwidth]{\detokenize{figures/videosearch/chatbotscreenshort videosearch.png}}
\caption{Video Search Portal — AI chatbot returning video results for a natural language query}
\label{fig:screen_vp_chat}
\end{figure}


\begin{figure}[H]
\centering
\includegraphics[width=\textwidth]{\detokenize{figures/ar/ar admin custom quotes.png}}
\caption{Admin — Pending Custom Quotes: review quote requests and assign a Stripe price ID}
\label{fig:admin_custom_quotes}
\end{figure}

\begin{figure}[H]
\centering
\includegraphics[width=\textwidth]{\detokenize{figures/ar/ar admin quotes status.png}}
\caption{Admin — Custom Quotes Status: track payment state and account creation for each quote}
\label{fig:admin_quotes_status}
\end{figure}

\begin{figure}[H]
\centering
\includegraphics[width=\textwidth]{\detokenize{figures/ar/ar admin submited video review.png}}
\caption{Admin — Submitted Videos: approve or reject user-submitted clips before library ingestion}
\label{fig:admin_submitted_videos}
\end{figure}

\begin{figure}[H]
\centering
\includegraphics[width=\textwidth]{\detokenize{figures/ar/ar admin user conversations.png}}
\caption{Admin — User Conversations: inspect full chat transcripts across all users}
\label{fig:admin_user_conversations}
\end{figure}

\begin{figure}[H]
\centering
\includegraphics[width=\textwidth]{\detokenize{figures/ar/ar admin chatbot conversationt tester.png}}
\caption{Admin — Chatbot Tester: test the assistant in chat and video-search modes with saved sessions}
\label{fig:admin_chatbot_tester}
\end{figure}

\begin{figure}[H]
\centering
\includegraphics[width=\textwidth]{\detokenize{figures/ar/ar admin faq manager.png}}
\caption{Admin — FAQ Manager: create, edit, and publish FAQ entries backed by Qdrant vector storage}
\label{fig:admin_faq_manager}
\end{figure}

\section*{Conclusion}
\addcontentsline{toc}{section}{Conclusion}

Sprint 5 closes the agentic layer and wires both portals to it. The final v2 design uses explicit
routing, deterministic video-search planning, the advanced hybrid search endpoint, Redis-capable
conversation memory, citation metadata for RAG, and guarded fallback paths. The production search
path is now faster, cheaper, and more testable. The Whisper pipeline lets users dictate queries,
and the rights portal exposes an admin control surface for moderation, conversation supervision,
chatbot testing, and vector-backed FAQ maintenance. The last sprint now focuses on testing,
deployment, and observability.
